\section{Experimental Protocol}
\label{sec:protocol}

\subsection{Map split and query selection}

The checksum-pinned source snapshot contains Moving AI pathfinding maps and
random scenario files \citep{sturtevant2012benchmarks}.  Table~\ref{tab:design}
summarizes the map-disjoint design.  The same four topology families occur on
both sides, but no map file overlaps.  Development uses one map per family and
40 queries per map to validate correctness, completeness, provenance, and the
analysis code; it performs no selection.  Evaluation begins only after an
external audit of those checks and uses two new maps per family with 100
queries per map.

\begin{table}[t]
\centering
\small
\caption{Experimental design.  Each query-method pair has one warm-up and eight
timed searches.}
\label{tab:design}
\begin{tabular}{lrrr}
\toprule
Split & Maps & Queries & Search runs \\
\midrule
Development & 4 & 160 & 5,760 \\
Sealed evaluation & 8 & 800 & 28,800 \\
\midrule
Total & 12 & 960 & 34,560 \\
\bottomrule
\end{tabular}
\end{table}

For each map, scenario files 1--4 are read in original source order.  A row is
valid when endpoints are distinct traversable cells in the same four-neighbor
component.  The plan takes the first 10 valid rows per file for development
and the first 25 for evaluation.  Duplicate $(\text{map},s,t)$ keys invalidate
the plan.  No query is sampled, replaced, or removed after timing.

The four development maps are
\texttt{maze-128-128-1}, \texttt{random-64-64-10},
\texttt{room-64-64-8}, and \texttt{warehouse-20-40-10-2-1}.
The evaluation's validation-source maps are
\texttt{maze-128-128-2}, \texttt{random-32-32-20},
\texttt{room-64-64-16}, and \texttt{warehouse-10-20-10-2-1};
its holdout-source maps are \texttt{maze-128-128-10},
\texttt{random-32-32-10}, \texttt{room-32-32-4}, and
\texttt{warehouse-10-20-10-2-2}.

\subsection{Balanced timing}

Within every query, the base method order is Manhattan, eager full, lazy full,
and staged.  One untimed warm-up uses that order.  Eight timed repetitions
left-rotate it by repetition index, cycling through the four positions twice.
Thus every method appears exactly twice in every position for every query.
Runs are sequential on one machine (Windows 11, CPython 3.12.9, 8 logical
CPUs); no concurrent search workers are used.

For a query and method, search time is the ordinary even median of eight
observations: the mean of the fourth and fifth sorted values.  Repetitions
reduce measurement noise and are not statistical units.  The primary
staged/lazy estimand is
\begin{equation}
 \exp\!\left(\operatorname{median}_{q}
     \log\frac{T_{q,\mathrm{staged}}}{T_{q,\mathrm{lazy}}}\right).
 \label{eq:ratio}
\end{equation}
Its percentile interval uses 10,000 deterministic bootstrap replicates
(seed 23725513), sampling maps first and then queries within sampled maps.
The eight maps are the top-level clusters.  Quartiles use Hyndman--Fan type 7.

\subsection{Validation and analysis}

All 800 records must be present, all costs must match replayed BFS, every path
must be legal, method summaries must match all repetitions, and the three
full-landmark expansion digests must agree.  Timing is analyzed only after all
checks pass.  There are no post-hoc exclusions.  Mechanism measures are exact per-query
counter differences.  Spearman associations use average ranks and are
descriptive.  Preprocessing amortization computes, within each map,
\begin{equation}
 A_m(Q)=\overline{T}_{m}+B_m/Q,
 \label{eq:amortization}
\end{equation}
where $B_m$ is measured landmark construction time and $Q$ is queries served;
the reported value equal-weights the eight maps.  Manhattan has $B_m=0$;
all landmark modes receive the same build charge.
